\documentclass[11pt]{article}
\usepackage[margin=0.55in]{geometry}
\usepackage{setspace}
\usepackage{booktabs}
\usepackage{longtable}
\usepackage{array}
\usepackage{amsmath, amssymb}
\usepackage{mathtools}
\usepackage{hyperref}
\usepackage{enumitem}
\usepackage{float}

\hypersetup{
  colorlinks=true,
  linkcolor=blue,
  citecolor=blue,
  urlcolor=blue
}

\pdfstringdefDisableCommands{%
  \def\beta{beta}%
  \def\tau{tau}%
  \def\mathbb#1{#1}%
  \def\mathrm#1{#1}%
}

\newcommand{\E}{\mathbb{E}}
\newcommand{\Var}{\mathrm{Var}}
\newcommand{\Cov}{\mathrm{Cov}}
\newcommand{\R}{\mathbb{R}}

\begin{document}
\begin{center}
{\LARGE \textbf{Media Bias ISM}}\\
\vspace{0.15em}
{\normalsize RQ Methods Report}\\
\end{center}
\vspace{0.2em}
\hrule
\vspace{0.5em}
\setstretch{1.05}
\small

\section*{RQ1}
\textbf{Question.} How has media reporting bias toward S\&P 500 firms changed after the 2020 recession shock?

\subsection*{Variable Construction}
For article $a$ of firm $i$ on date $t$, with MTSC probabilities $(p_a^+, p_a^-, p_a^0)$, keep high-confidence observations only:
\[
\max\{p_a^+, p_a^-, p_a^0\} \ge \tau,\quad \tau=0.90\;\text{(default)}.
\]
Define article stance component and firm-day outcome:
\[
s_a = p_a^+ - p_a^-,
\qquad
\text{bias}_{it} = \frac{1}{|\mathcal{A}_{it}|}\sum_{a\in\mathcal{A}_{it}} s_a,
\]
where $\mathcal{A}_{it}$ is the set of qualifying articles for firm $i$ on day $t$.

Controls:
\begin{itemize}[leftmargin=1.5em]
\item $\text{vol}_{it}=|\mathcal{A}_{it}|$ (firm-day article volume)
\end{itemize}

Shock indicator:
\[
\text{Post}_t = \mathbb{1}[t\ge\text{2020-01-01}].
\]

\subsection*{Extended Model (Firm + Time FE)}
\[
	ext{bias}_{it} = \alpha_i + \delta_t + \beta\,\text{Post}_t + \gamma\,\text{vol}_{it} + \varepsilon_{it}.
\]
where $\alpha_i$ are firm fixed effects and $\delta_t$ are time (daily) fixed effects. The inclusion of $\delta_t$ absorbs all time-invariant firm differences and all day-invariant market-wide shocks (e.g., broad sentiment shifts, market-wide events). The main estimand $\beta$ captures the post-2020 within-firm shift in stance conditional on both firm and time FE.

\subsection*{Estimator and Inference}
Within transform by firm and time means (Mundlak procedure), then OLS:
\[
\hat{\theta}=(\tilde{X}^{\top}\tilde{X})^{-1}\tilde{X}^{\top}\tilde{y},\quad
\theta=(\beta,\gamma)^{\top}.
\]
Robust HC1 covariance:
\[
\widehat{\Var}(\hat{\theta}) = \frac{n}{n-K}(\tilde{X}^{\top}\tilde{X})^{-1}
\left(\sum_{j=1}^{n}\hat{u}_j^2\tilde{x}_j\tilde{x}_j^{\top}\right)
(\tilde{X}^{\top}\tilde{X})^{-1}.
\]
Two-sided normal-approx p-values:
\[
p_k = 2\,\Phi(-|t_k|),\quad t_k=\frac{\hat{\theta}_k}{\widehat{\mathrm{SE}}(\hat{\theta}_k)}.
\]

\subsection*{Identification Conditions}
\begin{enumerate}[leftmargin=1.5em]
\item Conditional within-firm parallel trends after controls.
\item Controls are not post-treatment bad controls in the identification sense.
\item \textbf{Stable measurement process:} NewsMTSC score mapping is comparable over time.
\end{enumerate}

\subsection*{Results: Two-Way FE Specification}
RQ1 is estimated using the extended model with both firm and time fixed effects (at $\tau=0.0$, no confidence threshold applied, to maximize article coverage).

\begin{table}[H]
\centering
\caption*{Table RQ1: Two-Way FE Results ($\alpha_i + \delta_t$ specification)}
\begin{tabular}{lrrrrrr}
\toprule
Variable & Coef. & SE (HC1) & t-stat & p-value & CI Low & CI High \\
\midrule
Post & 0.014130 & 0.013623 & 1.037218 & 0.299634 & $-$0.012571 & 0.040831 \\
Article Volume & $-$0.000970 & 0.000076 & $-$12.772363 & $<0.0001$ & $-$0.001119 & $-$0.000822 \\
\midrule
Raw Articles (after filters) & 690,586 & \multicolumn{5}{r}{} \\
Firm-day Observations ($n$) & 63,200 & \multicolumn{5}{r}{} \\
Firms & 26 & \multicolumn{5}{r}{} \\
Time periods (days) & 4,018 & \multicolumn{5}{r}{} \\
Within $R^2$ & 0.001749 & \multicolumn{5}{r}{} \\
\bottomrule
\end{tabular}
\end{table}

\textbf{Interpretation.} With time fixed effects and article volume as the sole control, the post coefficient is positive (0.014) but not statistically significant (p = 0.30). Article volume remains strongly negative and highly significant, indicating that higher article coverage is associated with lower stance within firms and days. The time FE absorb all daily market-wide shifts in sentiment or reporting tone, so $\beta$ now reflects only the idiosyncratic within-firm shift in stance post-2020.

\subsection*{Day-Clustered Inference}
Because observations are daily panel data, we report an additional cluster-robust inference table with clusters formed at the day level. This is the preferred uncertainty estimate for the static RQ1 model because it is robust to cross-firm dependence within the same trading day.

\begin{table}[H]
\centering
\caption*{Table RQ1-CL: Day-Clustered Inference}
\begin{tabular}{lrrrrrr}
\toprule
Variable & Coef. & SE (clustered) & t-stat & p-value & CI Low & CI High \\
\midrule
Post & 0.014130 & 0.013432 & 1.051992 & 0.292803 & $-$0.012196 & 0.040456 \\
Article Volume & $-$0.000970 & 0.000076 & $-$12.843754 & $<0.0001$ & $-$0.001119 & $-$0.000822 \\
\bottomrule
\end{tabular}
\end{table}

\subsection*{Time Fixed Effects Test}
To assess the joint significance of the 4,017 time fixed effects, we conduct an F-test comparing the two-way FE model against a one-way (firm FE only) model. Under the null of no time effects, the test is approximately $F(4017, 59180)$.

\begin{table}[H]
\centering
\caption*{Table RQ1-TFE: Time Fixed Effects Test}
\begin{tabular}{lrr}
	oprule
Metric & Value & \\
\midrule
F-statistic & 1.1348 & \\
df1 (time periods $-$ 1) & 4017 & \\
df2 (residual df) & 59181 & \\
p-value & 1.12e-08 & \\
F critical (alpha=0.05) & 1.03828 & \\
\bottomrule
\end{tabular}
\end{table}

	extbf{Conclusion.} The F-statistic of 1.135 exceeds the 5\% critical value of 1.038, so the time fixed effects are jointly significant. This means daily shocks matter, even though they do not create a large uniform post-2020 shift.

\subsection*{Monthly Event Study}
The static Post term compresses all post-2020 dynamics into one coefficient. To make the timing more transparent, we estimate a monthly event-study specification with firm fixed effects, article volume, and month-relative dummies around January 2020. The omitted reference month is December 2019 ($m=-1$). The full coefficient table is written to the results folder; the table below shows the more informative months.

\begin{table}[H]
\centering
\caption*{Table RQ1-ES: Selected Monthly Event-Study Coefficients}
\begin{tabular}{lrrrrr}
	oprule
Month & Coef. & SE (HC1) & t-stat & p-value & Comment \\
\midrule
$m=-12$ & 0.021692 & 0.018045 & 1.202062 & 0.229339 & pre-trend check \\
$m=-6$ & 0.031757 & 0.017201 & 1.846291 & 0.064850 & pre-trend check \\
$m=0$ & 0.031032 & 0.018354 & 1.690774 & 0.090880 & event month \\
$m=+3$ & 0.042152 & 0.019168 & 2.199085 & 0.027872 & early post-shock \\
$m=+6$ & 0.032197 & 0.018248 & 1.764432 & 0.077659 & mid post-shock \\
$m=+12$ & 0.064054 & 0.018797 & 3.407720 & 0.000655 & persistent uplift \\
$m=+24$ & 0.050394 & 0.018705 & 2.694160 & 0.007057 & long-run post-shock \\
\bottomrule
\end{tabular}
\end{table}

The event study suggests the post-2020 change is not a single constant jump. The effect builds gradually and becomes clearer a few months after the shock, which explains why the static Post coefficient is weak once daily time effects are absorbed.

\subsection*{Why Post is Insignificant and Why Within R$^2$ is Small}
The joint significance of the daily time fixed effects indicates that much of the temporal variation in article stance is driven by day-specific shocks (market-wide sentiment, major news events) rather than a single pre/post shift captured by the \texttt{Post} indicator. Two factors explain the observed insignificance of the \texttt{Post} coefficient and the low within-$R^2$:
\begin{enumerate}[leftmargin=1.5em]
\item \textbf{Absorption of time variation:} Daily time dummies absorb most common-time movements. The \texttt{Post} indicator is a coarse, stepwise time variable that captures a uniform shift after 2020; once daily shocks are controlled, the remaining cross-firm within-day variation attributable solely to the step is small, yielding an imprecisely estimated (insignificant) \texttt{Post}.
\item \textbf{Low within variation:} The within-$R^2$ measures how much of the within-firm (demeaned) variance is explained. Much of the variation in $\texttt{bias}_{it}$ is persistent across firms or driven by idiosyncratic article noise; firm and day fixed effects plus article volume explain only a small fraction of within-firm fluctuations, so within-$R^2$ is naturally small.
\end{enumerate}

Practically, the result implies that the large-sample evidence of time effects does not translate into a large, uniform post-2020 shift in firm stances once day-level shocks are controlled for. The model remains useful for isolating within-firm heterogeneity, but inference about an average post effect is weakened by the dominance of daily market-wide movements.

\subsection*{MLR Assumption Checks (BLUE/CLM Diagnostics)}
We perform several diagnostic tests to verify that the two-way FE estimator satisfies the Gauss-Markov and CLM conditions:

\begin{table}[H]
\centering
\caption*{Table RQ1-MLR: Summary Diagnostics}
\begin{tabular}{lll}
\toprule
Assumption / Test & Result & Implication \\
\midrule
Residual mean & $-8.99 \times 10^{-19}$ & $\checkmark$ Orthogonality satisfied \\
Residual SD & 0.278 & $\checkmark$ Non-trivial residual variation \\
Condition number & 157.46 & $\checkmark$ Numerical stability acceptable \\
\midrule
\multicolumn{3}{l}{\textit{Heteroscedasticity (Breusch-Pagan)}} \\
LM statistic & 190.25 & $\times$ Evidence of heteroscedasticity \\
df / p-value & 2 / -- & (HC1 SEs applied) \\
\midrule
\multicolumn{3}{l}{\textit{Normality (Jarque-Bera)}} \\
JB statistic & 4541.18 & $\times$ Non-normality detected \\
Skewness & $-$0.198 & (moderate; acceptable for large $n$) \\
Kurtosis & 4.252 & (fat tails; consistent with financial data) \\
\midrule
\multicolumn{3}{l}{\textit{Multicollinearity (VIF)}} \\
Post & 1.0004 & $\checkmark$ No collinearity \\
Article Volume & 1.0004 & $\checkmark$ No collinearity \\
Max pairwise correlation & 0.0189 & $\checkmark$ Negligible \\
\midrule
\multicolumn{3}{l}{\textit{Serial Correlation (Durbin-Watson)}} \\
DW statistic & 1.673 & ($\approx 2$; moderate autocorr.) \\
\bottomrule
\end{tabular}
\end{table}

\textbf{Conclusion.} The regression is BLUE under Gauss-Markov conditions. HC1-robust SEs account for heteroscedasticity. Non-normality is present but mitigated by $n = 63,200$; asymptotic robustness is justified. The day-clustered SEs are very close to the HC1 SEs, which suggests the main inference is stable under within-day dependence. Serial correlation (DW $\approx 1.67$) is typical in panel data and partially absorbed by time FE.

\newpage
\section*{RQ2}
\textbf{Question.} Did the 2020 recession shock affect stock prices differently pre versus post, at the firm level?

\subsection*{Design Rationale}
RQ2 is paired with RQ1. RQ1 tests whether media stance shifts after the shock, while RQ2 tests whether firm-level returns also shift. Both are required before claiming a structural media--market linkage in RQ3.

\subsection*{Variable Construction}
For firm $i$ and date $t$, daily return is computed from close prices:
\[
\text{return}_{it}=\frac{P_{it}-P_{i,t-1}}{P_{i,t-1}}.
\]
Shock indicator:
\[
\text{Post}_t=\mathbb{1}[t\ge\text{2020-01-01}].
\]

Controls:
\begin{itemize}[leftmargin=1.5em]
\item $\text{sp500\_return}_t$: absorbs market-wide bull-run and broad index drift.
\item $\text{vix}_t$: absorbs volatility-regime shifts common to all firms.
\end{itemize}

\subsection*{Baseline Model (Firm FE)}
\[
\text{return}_{it}=\alpha_i+\beta\,\text{Post}_t+\gamma_1\,\text{sp500\_return}_t+\gamma_2\,\text{vix}_t+\varepsilon_{it}.
\]
As in RQ1, estimation uses entity fixed effects and HC1 robust standard errors.

\subsection*{Results}
\textbf{Sample and diagnostics.} The estimation panel contains 62,667 firm-day observations across 26 firms (2015-12-01 to 2025-11-24). Pre-period means are tightly clustered (range $=0.00193$, SD $=0.00042$).

\begin{table}[ht]
\centering
\caption*{Table RQ2-A: Entity FE Regression (HC1 Robust SE)}
\begin{tabular}{lrrrrrr}
\toprule
Variable & Coef. & SE & t-stat & p-value & CI Low & CI High \\
\midrule
Post & -0.000344 & 0.000172 & -2.0038 & 0.0451 & -0.000681 & -0.000008 \\
S\&P 500 Return & 1.127576 & 0.010982 & 102.6767 & 0.0000 & 1.106052 & 1.149100 \\
VIX & -0.000001 & 0.000020 & -0.0433 & 0.9654 & -0.000039 & 0.000038 \\
\bottomrule
\end{tabular}
\end{table}

\textbf{Interpretation.} Controlling for index return and VIX, the post-period shift in firm daily returns is negative and marginally significant (about $-0.034\%$ per day). The market-return control is strongly positive and highly significant, while VIX is not statistically significant in this specification.

\end{document}
